\subsection{CT7 --- Exchange trichotomy}\label{ct:7}\label{the-exchange-trichotomy}
\ctsignature{\(\begin{array}{@{}l@{}}P_{\ast\to7}\to \ctcert{C1,C2,C3};\\ \ctint{P_{7\to3},P_{7\to12}};\ \ctrec{P_{7\to10},P_{7\to16}}.\end{array}\)}

CT7 is an exact machine for a bounded same-interface exchange.  Its stable
entry interface contains the ambient and branch states, the two local
representatives, their common boundary, the exchange witness, the context
system, the finite-scope predicate, and the C1--C3 claims.  The executable
framework adds context completeness, realization, distinction, and neutrality
predicates together with exact CT3, CT10, CT12, and CT16 entry contracts.

The core plan has six fields: \texttt{scope}, \texttt{context},
\texttt{realization}, \texttt{distinction}, \texttt{defect}, and
\texttt{neutral}.  The context node has one successor and constructs a
\texttt{ContextState} containing the completeness proof for the finite
generator.  The realization node either returns a C1 certificate or an
\texttt{UnrealizedState}.  Only the latter reaches the distinction node,
which returns a \texttt{DefectState} or a \texttt{NeutralState}.

Defect routing and neutral compression are separate APIs.  The defect node
returns C3, an exact CT3 input, or an exact CT12 peeling input.  The neutral node
returns C2, an exact CT10 input, or an exact CT16 whole-object input.  Thus a node
implementer proves only one local classification and may rely on every earlier
certificate through the predecessor index.  No string tag or unindexed
residual can bypass these states.

\begin{itemize}
\tightlist
\item \textbf{S-Def} is fixed by the validated exchange entry.
\item \textbf{S-Equiv} is the one-successor context-completeness certificate and is retained by both later branches.
\item \textbf{S-Dich} is discharged separately by realization and distinction.
\item \textbf{S-Comp} is discharged only at the neutral node, where a C2 certificate proves the smaller equivalent representative.
\item \textbf{S-Rout} and \textbf{S-Trig} are discharged by the four indexed handoff families and the consumer plan.
\end{itemize}

No finite context generator yields the typed scope terminal.  An assertion of
neutrality without generator completeness is not an executable branch: it is
an uninhabited node contract that must be repaired before admission.

\begin{figure}[p]
\centering
\resizebox{0.98\textwidth}{!}{%
\begin{tikzpicture}[x=1cm,y=1cm,every node/.style={font=\scriptsize}]
\node[bcwide] (entry) at (0,0) {{\tiny\ttfamily CT7.entry}\\\textbf{Validated exchange pair}};
\node[bcdec,bclemma,text width=3.45cm,minimum width=4.20cm] (scope) at (0,-2.15) {{\tiny\ttfamily CT7.decide.scope}\\\textbf{Finite complete generator?}};
\node[bcscopefail,text width=3.10cm,minimum width=3.10cm] (scopeout) at (7.35,-2.15) {{\tiny\ttfamily CT7.terminal.scope}\\\textbf{Scope exit}};
\node[bcwide,bclemma] (context) at (0,-4.35) {{\tiny\ttfamily CT7.certify.context}\\\textbf{S-Equiv completeness certificate}};
\node[bcdec,bclemma,text width=3.45cm,minimum width=4.20cm] (realize) at (0,-6.55) {{\tiny\ttfamily CT7.decide.realization}\\\textbf{Compatible target realization?}};
\node[bcclosed] (c1) at (7.35,-6.55) {{\tiny\ttfamily CT7.terminal.c1}\\\textbf{C1}};
\node[bcdec,bclemma,text width=3.45cm,minimum width=4.20cm] (distinction) at (0,-8.95) {{\tiny\ttfamily CT7.decide.distinction}\\\textbf{Distinguishing or neutral?}};
\node[bcroutechoice,bctheorem,text width=3.25cm,minimum width=4.10cm] (defect) at (-5.30,-11.75) {{\tiny\ttfamily CT7.decide.defect}\\\textbf{Route distinguishing defect}};
\node[bcroutechoice,bctheorem,text width=3.25cm,minimum width=4.10cm] (neutral) at (5.30,-11.75) {{\tiny\ttfamily CT7.decide.neutral}\\\textbf{Compress or route neutral pair}};
\node[bcclosed,text width=2.80cm,minimum width=2.80cm] (c3) at (-10.10,-15.00) {{\tiny\ttfamily CT7.terminal.c3}\\\textbf{C3}};
\node[bcroute,bcrouteneutral,bcoutputroute,text width=2.95cm,minimum width=2.95cm] (ct3) at (-5.30,-15.00) {{\tiny\ttfamily CT7.terminal.ct3}\\\textbf{CT 3 exit}};
\node[bcroute,bcrouteneutral,bcoutputroute,text width=2.95cm,minimum width=2.95cm] (ct12) at (-0.50,-15.00) {{\tiny\ttfamily CT7.terminal.ct12}\\\textbf{CT 12 exit}};
\node[bcclosed,text width=2.80cm,minimum width=2.80cm] (c2) at (3.65,-18.15) {{\tiny\ttfamily CT7.terminal.c2}\\\textbf{C2}};
\node[bcroute,bcrouteneutral,bcdesign,text width=2.95cm,minimum width=2.95cm] (ct10) at (8.15,-18.15) {{\tiny\ttfamily CT7.terminal.ct10}\\\textbf{CT 10 exit}};
\node[bcroute,bcrouteneutral,bcdesign,text width=2.95cm,minimum width=2.95cm] (ct16) at (12.65,-18.15) {{\tiny\ttfamily CT7.terminal.ct16}\\\textbf{CT 16 exit}};
\draw[bcarr] (entry)--(scope);
\draw[bcscopearr] (scope)--node[bclabel,above]{exit}(scopeout);
\draw[bcarr] (scope)--node[bclabel,right]{ready}(context);
\draw[bcarr] (context)--node[bclabel,right]{certified}(realize);
\draw[bcarr] (realize)--node[bclabel,above]{realized}(c1);
\draw[bcarr] (realize)--node[bclabel,right]{unrealized}(distinction);
\draw[bcarr] (distinction)--node[bclabel,above,sloped]{defect}(defect);
\draw[bcarr] (distinction)--node[bclabel,above,sloped]{neutral}(neutral);
\draw[bcarr] (defect)--node[bclabel,above,sloped]{routed}(c3);
\draw[bchandoff] (defect)--node[bclabel,right]{\(P_{7\to3}\)}(ct3);
\draw[bchandoff] (defect)--node[bclabel,above,sloped]{\(P_{7\to12}\)}(ct12);
\draw[bcarr] (neutral)--node[bclabel,above,sloped]{compress}(c2);
\draw[bcrecoverarr] (neutral)--node[bclabel,right]{\(P_{7\to10}\)}(ct10);
\draw[bcrecoverarr] (neutral)--node[bclabel,above,sloped]{\(P_{7\to16}\)}(ct16);
\end{tikzpicture}%
}
\caption[Closure-tactic 7: exact exchange machine]{Closure-tactic 7 as an exact sequential machine.  Context completeness is certified once.  Realization, distinction, defect routing, and neutral compression are independent node contracts.  Each edge carries the exact indexed state consumed by its successor.}
\label{fig:ct7-exchange-trichotomy}
\end{figure}
\FloatBarrier
